
\documentclass[border=4pt]{standalone}
\usepackage{polyglossia}
\setdefaultlanguage{russian}
\setotherlanguage{english}
\usepackage{fontspec}
\IfFontExistsTF{CMU Serif}{\setmainfont{CMU Serif}}{\setmainfont{Liberation Serif}}
\IfFontExistsTF{CMU Sans Serif}{\setsansfont{CMU Sans Serif}}{\setsansfont{Liberation Sans}}
\IfFontExistsTF{CMU Typewriter Text}{\setmonofont{CMU Typewriter Text}}{\setmonofont{DejaVu Sans Mono}}
\usepackage{amsmath,amssymb,mathtools,bm}
\usepackage{xcolor}
\definecolor{accent}{HTML}{1F4E79}
\definecolor{accent2}{HTML}{B85450}
\definecolor{shade}{HTML}{F4F1EA}
\colorlet{prosvBlue}{accent}
\colorlet{prosvRed}{accent2}
\colorlet{prosvLight}{accent!12}
\colorlet{prosvCream}{shade}
\definecolor{prosvGold}{HTML}{E8B647}
\definecolor{prosvGreen}{HTML}{6A9F58}
\colorlet{prosvGray}{black!55}
\usepackage{tikz}
\usetikzlibrary{calc, arrows.meta, decorations.pathreplacing,
                positioning, shapes.geometric, shapes.misc, fit,
                backgrounds}
\usepackage{pgfplots}
\pgfplotsset{compat=1.18}
\newcommand{\R}{\mathbb{R}}
\newcommand{\N}{\mathbb{N}}
\newcommand{\Z}{\mathbb{Z}}
\newcommand{\eps}{\varepsilon}
\newcommand{\norm}[1]{\left\|#1\right\|}
\newcommand{\abs}[1]{\left|#1\right|}
\newcommand{\NOD}{\gcd}
\newcommand{\Eul}{\varphi}
\newcommand{\Zn}[1]{\mathbb{Z}_{#1}}
\newcommand{\modn}[1]{\,(\mathrm{mod}\,#1)}
\newcommand{\term}[1]{\textcolor{accent2}{\textbf{#1}}}
\newcommand{\engterm}[1]{\textit{#1}}
\DeclareMathOperator*{\argmin}{arg\,min}
\DeclareMathOperator*{\argmax}{arg\,max}
\begin{document}

\begin{tikzpicture}[
  font=\small, every node/.style={align=center},
  circ/.style={circle, draw, minimum size=0.7cm}
]
  
  \node[circ, fill=yellow!50!brown!20] (g) at (0,3) {};
  \node[below, font=\scriptsize] at (0,2.6) {общая «жёлтая»};
  
  \node[circ, fill=red!60!brown!50] (sa) at (-5,3) {};
  \node[below, font=\scriptsize, align=center] at (-5,2.6) {Алисин «красный»\\(секрет)};
  
  \node[circ, fill=blue!60!cyan!50] (sb) at (5,3) {};
  \node[below, font=\scriptsize, align=center] at (5,2.6) {Бобов «синий»\\(секрет)};
  
  
  \node[circ, fill=orange!70!red!30] (ya) at (-2.5,1.0) {};
  \node[below, font=\scriptsize, align=center] at (-2.5,0.6) {жёлт. $+$ красн.\\(передаётся\\открыто)};
  
  \node[circ, fill=green!60!blue!30] (yb) at (2.5,1.0) {};
  \node[below, font=\scriptsize, align=center] at (2.5,0.6) {жёлт. $+$ син.\\(передаётся\\открыто)};
  
  \draw[->, prosvBlue] (g) -- (ya);
  \draw[->, prosvBlue] (sa) -- (ya);
  \draw[->, prosvBlue] (g) -- (yb);
  \draw[->, prosvBlue] (sb) -- (yb);
  
  
  \node[circ, fill=brown!50] (final) at (0,-2.0) {};
  \node[below, font=\scriptsize, align=center] at (0,-2.4) {итоговая «коричневая»\\ --- общий секрет};
  
  \draw[->, prosvGreen, thick] (yb) -- node[above right, font=\scriptsize, text=prosvGreen]{+ Алисин «красный»} (final);
  \draw[->, prosvGreen, thick] (ya) -- node[above left, font=\scriptsize, text=prosvGreen]{+ Бобов «синий»} (final);
\end{tikzpicture}
\end{document}
